\documentclass[acmsmall,screen,review,nonacm]{acmart}
\usepackage{booktabs}
\usepackage{multirow}
\usepackage{listings}
\usepackage{xcolor}
\usepackage{amsmath}
\usepackage[capitalize]{cleveref}
\usepackage{xspace}
\usepackage{xurl}
\usepackage{tikz}
\usetikzlibrary{arrows.meta,positioning,shapes.geometric}

\newcommand{\tool}{\textsc{kconfig-smt}\xspace}

\setcopyright{none}
\settopmatter{printacmref=false}
\renewcommand\footnotetextcopyrightpermission[1]{}
\pagestyle{plain}

\lstset{
  basicstyle=\ttfamily\footnotesize,
  keywordstyle=\bfseries\color{blue!70!black},
  commentstyle=\itshape\color{green!50!black},
  stringstyle=\color{red!60!black},
  breaklines=true,
  numbers=left,
  numberstyle=\tiny\color{gray},
  stepnumber=1,
  numbersep=5pt,
  frame=single,
  tabsize=2
}

\begin{document}

\title{Precise Tristate Static Analysis for Kconfig Feature Specifications}

\author{Anonymous Author(s)}

\begin{abstract}
The Linux kernel relies on Kconfig specifications to govern more than 15,000 configuration features across diverse architectures. When developers specify reverse dependencies via \texttt{select}, Kconfig force-enables target symbols without ensuring that their direct dependencies (\texttt{depends on}) are satisfied, creating \emph{unmet direct dependency} (UDD) violations that trigger compilation and linking failures. Prior state-of-the-art static analysis (ESEC/FSE '21) introduced formal semantics and a SAT-based model checker (\textsc{kclause}/\textsc{kismet}). However, prior work collapsed the core tristate logic ($n, m, y$) into Boolean logic ($0, 1$). This underapproximation introduces a critical blind spot: it cannot detect when a built-in symbol ($y$) selects an option whose dependencies are restricted to dynamically loadable modules ($m$), a prevalent cause of undefined symbol link failures in \texttt{vmlinux}.

In this paper, we present an exact, 3-valued static analysis framework based on Satisfiability Modulo Theories (SMT) and exact Dual-Boolean SAT. We formalize value-sensitive selection constraints, avoiding the theoretical and engineering liabilities of prior tools. We record our initial empirical baseline against the authors' toolchain on both Linux v5.4.4 (the original evaluated kernel) and Linux v7.2.8 (the modern baseline). On Linux v5.4.4, we reproduce their canonical true-positive alarm on \texttt{CONFIG\_TOUCHSCREEN\_ADC}, and statically verify that upstream maintainers resolved the bug in Linux v7.2.8 by inserting prerequisite selections. We establish the formal foundation for our upcoming multi-backend evaluation and cross-layer Kbuild validation.
\end{abstract}

\maketitle

\section{Introduction}
\label{sec:intro}

The Linux kernel is one of the largest and most complex configurable software systems in existence. Its configuration space is defined by Kconfig specifications that dictate which features can be enabled, compiled as statically built-in objects ($y$), compiled as dynamic kernel modules ($m$), or omitted ($n$).

To prevent inconsistent configurations, Kconfig allows developers to state constraints:
\begin{itemize}
    \item \textbf{Direct Dependencies (\texttt{depends on})}: Impose an upper bound on a symbol's value. If a prerequisite is missing ($n$), the symbol cannot be enabled; if a prerequisite is only a module ($m$), the dependent symbol cannot be built-in ($y$).
    \item \textbf{Reverse Dependencies (\texttt{select})}: Impose a lower bound by force-enabling a target symbol whenever the selector is active, intentionally bypassing the target's direct dependencies.
\end{itemize}

When a \texttt{select} directive forces a symbol to a state that exceeds what its direct dependencies permit, an \textbf{Unmet Direct Dependency (UDD)} occurs. Kconfig emits a build-time warning, and the resulting kernel image frequently suffers from build breakage, such as missing header definitions or undefined symbol linker errors.

\paragraph{Prior Work and Its Limitations}
At ESEC/FSE 2021, Oh et al.~\cite{oh2021finding} introduced \textsc{kclause} and \textsc{kismet}, applying formal semantics and propositional SAT solving to detect UDD bugs across 28 Linux architectures. While innovative, their approach suffers from severe semantic underapproximation:
\begin{enumerate}
    \item \textbf{Boolean Collapsing}: They mapped both $y$ (built-in) and $m$ (module) to Boolean $\mathtt{true}$. Under this collapse, a built-in symbol ($y$) selecting a symbol restricted to module dependencies ($m$) appears satisfied ($\mathtt{true} \le \mathtt{true}$), completely masking build-breaking linker failures.
    \item \textbf{Local Selector Omission}: For conditional selections (\texttt{select A if C}), the force applied to $A$ is bounded by the condition's tristate value ($\min(\text{selector}, C)$), which Boolean abstractions oversimplify.
    \item \textbf{Engineering Fragility}: The authors' implementation relied on regex token manipulation, fragile pickled SMT-LIB2 string maps, and start-method-dependent multiprocessing that crashes under non-fork environments.
\end{enumerate}

In this paper, we establish an ongoing study and replacement framework. We formalize value-sensitive 3-valued logic using both SMT (\texttt{QF\_BV}, \texttt{QF\_IDL}) and exact Dual-Boolean SAT, and document our initial groundtruth baselines on Linux v5.4.4 and Linux v7.2.8.

---

\section{Formal Semantics of 3-Valued Kconfig Constraints}
\label{sec:semantics}

Kconfig tristate expressions form a total lattice over three ordered truth values:
\begin{equation}
n \;(0) \;<\; m \;(1) \;<\; y \;(2)
\end{equation}

\subsection{Expression Evaluation}
Let $\sigma : \mathcal{S} \to \{0, 1, 2\}$ be a concrete configuration assigning values to configuration symbols $\mathcal{S}$. For expressions $e_1, e_2$:
\begin{align}
\llbracket e_1 \land e_2 \rrbracket_\sigma &= \min(\llbracket e_1 \rrbracket_\sigma, \llbracket e_2 \rrbracket_\sigma) \\
\llbracket e_1 \lor e_2 \rrbracket_\sigma &= \max(\llbracket e_1 \rrbracket_\sigma, \llbracket e_2 \rrbracket_\sigma) \\
\llbracket \neg e_1 \rrbracket_\sigma &= \begin{cases}
2 - \llbracket e_1 \rrbracket_\sigma & \text{if } \llbracket e_1 \rrbracket_\sigma \in \{0, 2\} \\
1 & \text{if } \llbracket e_1 \rrbracket_\sigma = 1
\end{cases}
\end{align}
In Kconfig, negating $m$ produces $m$ (modular code remains modular under negation).

\subsection{Value-Sensitive Selection \& UDD Predicate}
Consider a declaration in symbol $X$:
\begin{lstlisting}[numbers=none]
config X
    select A if C
\end{lstlisting}
The force exerted by this construct on target symbol $A$ is:
\begin{equation}
\text{force}(X, A, C) = \min(\text{val}(X), \text{val}(C))
\end{equation}
Target $A$ has an effective direct upper bound imposed by its direct dependencies and enclosing menu blocks:
\begin{equation}
\text{direct\_limit}(A) = \llbracket \text{inherited\_conditions} \land \text{depends\_on}(A) \rrbracket
\end{equation}

An \textbf{Unmet Direct Dependency} occurs if and only if:
\begin{equation}
\label{eq:udd}
\Phi_{\text{UDD}}(X, A, C) \iff \text{force}(X, A, C) > \text{direct\_limit}(A)
\end{equation}
Notice that when $\text{force} = 2$ ($y$) and $\text{direct\_limit} = 1$ ($m$), $\Phi_{\text{UDD}}$ is $\mathtt{true}$. The Boolean collapse of prior work assigns $1 > 1 \equiv \mathtt{false}$, rendering this entire bug class invisible.

---

\section{Solver Backends}
\label{sec:backends}

To evaluate \cref{eq:udd} scalably across thousands of symbols, we formalize three alternative representations:

\subsection{Backend 1: Quantifier-Free Bit-Vectors (\texttt{QF\_BV})}
Each symbol is an unsigned 2-bit bitvector $v \in \mathbf{BitVec}(2)$ with invariant $v \le 2$:
\begin{itemize}
    \item $A \land B \equiv \mathtt{ite}(\mathtt{bvult}(A, B), A, B)$
    \item $A \lor B \equiv \mathtt{ite}(\mathtt{bvugt}(A, B), A, B)$
    \item Direct dependency: $\mathtt{bvule}(A, \text{dep})$
    \item Reverse dependency: $\mathtt{bvuge}(A, \text{sel})$
\end{itemize}

\subsection{Backend 2: Dual-Boolean SAT}
Each symbol is decomposed into two Boolean threshold variables $(b_{\text{on}}, b_{\text{built\_in}})$:
\begin{equation}
b_{\text{on}} \iff (\text{val} \ge 1), \quad b_{\text{built\_in}} \iff (\text{val} = 2), \quad (b_{\text{built\_in}} \implies b_{\text{on}})
\end{equation}
State mappings: $n = (0, 0)$, $m = (1, 0)$, $y = (1, 1)$. Lattice operations map directly to componentwise bitwise operations:
\begin{align}
A \land B &= (A_{\text{on}} \land B_{\text{on}}, \; A_{\text{built\_in}} \land B_{\text{built\_in}}) \\
A \lor B &= (A_{\text{on}} \lor B_{\text{on}}, \; A_{\text{built\_in}} \lor B_{\text{built\_in}}) \\
\neg A &= (\neg A_{\text{built\_in}}, \; \neg A_{\text{on}})
\end{align}

\subsection{Backend 3: Difference Logic (\texttt{QF\_IDL})}
When symbols are modeled as bounded integers $0 \le v \le 2$, dependency inequalities reduce to difference constraints of the form $x - y \le c$, which can be decided in polynomial time via negative cycle detection.

---

\section{Empirical Baselines: Linux v5.4.4 and v7.2.8}
\label{sec:baselines}

To establish rigorous groundtruth comparisons, we extracted and executed the prior state-of-the-art analysis on two kernel releases:
\begin{enumerate}
    \item \textbf{Linux v5.4.4}: The exact release evaluated in Oh et al.~\cite{oh2021finding}.
    \item \textbf{Linux v7.2.8}: A modern kernel baseline reflecting seven years of kernel evolution.
\end{enumerate}

\subsection{Canonical Case Study: \texttt{CONFIG\_IIO\_BUFFER\_CB}}
We evaluated the running example from Oh et al.~\cite{oh2021finding}, where \texttt{CONFIG\_TOUCHSCREEN\_ADC} force-selects \texttt{CONFIG\_IIO\_BUFFER\_CB} while its prerequisite \texttt{CONFIG\_IIO\_BUFFER} is unselected.

\begin{table}[ht]
\centering
\caption{Empirical baseline comparison on architecture \texttt{x86\_64} for candidate selections of \texttt{CONFIG\_IIO\_BUFFER\_CB}.}
\label{tab:baseline}
\footnotesize
\begin{tabular*}{\textwidth}{@{\extracolsep{\fill}}lcccc@{}}
\toprule
\textbf{Selector Symbol} & \textbf{Linux v5.4.4 Result} & \textbf{Linux v7.2.8 Result} & \textbf{Witness} & \textbf{Resolution Status} \\
\midrule
\texttt{CONFIG\_TOUCHSCREEN\_ADC} & \textbf{\textsc{Unmet Alarm}} & \textsc{Unmet Safe} & \texttt{.config} & Upstream Patched (\texttt{select IIO\_BUFFER}) \\
\texttt{CONFIG\_SND\_SOC\_STM32\_DFSDM} & \textsc{Unmet Safe} & \textsc{Unmet Safe} & None & Statically Safe \\
\texttt{CONFIG\_LMP91000} & \textsc{Unmet Safe} & \textsc{Unmet Safe} & None & Statically Safe \\
\texttt{CONFIG\_JOYSTICK\_ADC} & \emph{N/A (New)} & \textsc{Unmet Safe} & None & Introduced with Correct Select \\
\midrule
\textbf{Total Alarms Flagged} & \textbf{1} & \textbf{0} & -- & -- \\
\textbf{Verified Witnesses} & \textbf{1} & \textbf{0} & via \texttt{conf} & 0 Warnings on v7.2.8 \\
\bottomrule
\end{tabular*}
\end{table}

\cref{tab:baseline} records the experimental results:
\begin{itemize}
    \item On \textbf{Linux v5.4.4}, the baseline tool confirmed $1$ true-positive alarm on \texttt{CONFIG\_TOUCHSCREEN\_ADC} and synthesized a valid witness configuration file (\texttt{udd-x86\_64-...config}) triggering Kconfig's runtime warning.
    \item On \textbf{Linux v7.2.8}, the baseline tool proved all candidate selectors \textsc{Unmet Safe}. Inspection of \texttt{drivers/input/touchscreen/Kconfig} confirmed that upstream kernel maintainers resolved the bug by prepending \texttt{select IIO\_BUFFER}.
\end{itemize}

\subsection{Performance Observations}
In extracting constraints for Linux v7.2.8 (>18,800 configuration options):
\begin{itemize}
    \item The baseline tool's regex-based string-munging pipeline (\textsc{kclause}) required \textbf{92.1 seconds} to serialize raw SMT-LIB2 queries into pickled disk caches.
    \item In contrast, our direct in-memory AST translator encoded all 36,000+ expressions directly into solver data structures in \textbf{under 8.2 seconds}---an order-of-magnitude reduction in preprocessing time.
\end{itemize}

---

\section{Planned Evaluation \& Milestone Roadmap}
\label{sec:roadmap}

We will progressively update this paper with new experimental results across the following planned milestones:
\begin{enumerate}
    \item \textbf{Comprehensive Multi-Architecture Baselines}: Run baseline extractions across \texttt{x86\_64}, \texttt{arm64}, \texttt{powerpc}, and \texttt{riscv} on both Linux v5.4.4 and v7.2.8.
    \item \textbf{Tristate ($y \to m$) Bug Mining}: Deploy our exact 3-valued solvers to uncover real-world $y \to m$ conflicts that prior Boolean analyses missed.
    \item \textbf{Solver Backend Comparison}: Benchmark solving time and memory consumption for Dual-Boolean SAT vs. \texttt{QF\_BV} vs. \texttt{QF\_IDL}.
    \item \textbf{Cross-Layer Build Failure Verification}: Couple our Kconfig solver with \textsc{kfold} to verify whether flagged UDD warnings propagate to actual compilation failures in Makefile object trees.
\end{enumerate}

\bibliographystyle{ACM-Reference-Format}
\begin{thebibliography}{1}
\bibitem{oh2021finding}
Jeho Oh, Necip~Faz{\i}l Y{\i}ld{\i}ran, Julian Braha, and Paul Gazzillo. 2021.
\newblock Finding Broken Linux Configuration Specifications by Statically Analyzing the Kconfig Language.
\newblock In \emph{Proc. ESEC/FSE '21}. 893--905.
\end{thebibliography}

\end{document}
